\chapter{Chemical Equations: A Reaction Ledger}

\begin{kaichang}
\textbf{Translator safety note:} The concentrated-acid, sealed-flame, and hydrogen/oxygen examples are not home activities. Do not improvise them or ignite gas mixtures; a trained teacher must risk-assess any live demonstration. See \href{https://institute.acs.org/acs-center/lab-safety/education-training/safer-experiments.html}{ACS demonstration guidance}.

Light a birthday candle. It grows shorter until only a little wax remains. Where did the wax go? ``It burned,'' people say. But does burning mean disappearing into nothing?

Now leave an iron nail on a damp windowsill. Weeks later, red-brown rust covers it. Weigh before and after, and something strange appears: the rusty nail is \textbf{heavier}.

One thing disappears, another gains weight. Does matter obey rules? This chapter establishes a reaction ledger with a familiar name: the chemical equation.
\end{kaichang}

\section{Where Did the Wax Go? Weigh First}

Before answering, act like detectives and do only what can be measured.

Weigh a new candle, light it, let half burn, then weigh again. Unsurprisingly, it is lighter. Dry ice in an open dish shrinks and vanishes; concentrated hydrochloric acid left open also leaves its bottle lighter. Matter seems to escape.

But wait: these experiments all use \textbf{open} containers, admitting and releasing gases. What if we weighed not just the obvious reactants and products, but everything in the room?

Such false disappearances abound. A car empties its tank, yet every \ud{1}{L} of gasoline burned releases over \ud{2}{kg} of carbon dioxide. The car gets lighter, the air heavier. Spoiling food smells because organic molecules recombine with oxygen and microbial substances, releasing some products into air. Cut apples brown through quiet reactions with oxygen. Apparent mass creation is similar: a sapling's growth comes mostly not from eating soil, but from carbon dioxide in air and water. Faulty weighing often means faulty boundaries.

\begin{shiyan}
\textbf{A Bubbling Reaction in a Closed Bottle}

Materials: food-grade baking soda (sodium bicarbonate), white vinegar, a plastic bottle with a cap (or conical flask), a balloon, and a kitchen scale, preferably resolving \ud{0.1}{g}.

Procedure:
\begin{enumerate}[itemsep=1pt]
\item Pour about \ud{50}{mL} of vinegar into the bottle and put two small spoonfuls of baking soda in the balloon.
\item Fit the balloon over the bottle mouth without dropping the soda in. Weigh the entire apparatus, including all contents, and record its mass.
\item Lift the balloon so the soda falls into the vinegar. Bubbles appear immediately and the balloon inflates.
\item After the reaction settles, weigh everything again and compare with step 2.
\end{enumerate}

What to observe: inflation shows gas formation, yet total mass is almost unchanged, usually within scale uncertainty.

Safety note: these kitchen materials permit a home experiment, but work in ventilation, wash your hands afterward, and do not drink the liquid. Do not fit the balloon too tightly, to avoid bursting. Watch teachers or videos for more vigorous closed-container reactions such as combustion; do not attempt those yourself.
\end{shiyan}

Bubbles form and inflate the balloon, yet total mass stays put. No matter appeared or vanished: ingredients in soda and vinegar changed form, some becoming trapped carbon dioxide. Open the vessel and escaping gas makes the scale ``lie'' about lost mass. Close the doors and the books balance.

Return to wax: combustion mainly makes invisible carbon dioxide and water vapor that disperse. Weigh a candle sealed with its required air under a large glass cover and total mass remains conserved. A rusty nail gains oxygen from air, and that oxygen's atoms account for the extra mass (Chapter 15).

\section{Atoms Do Not Disappear; They Change Partners}

The observation is clear: \textbf{total mass stays unchanged through a reaction}. Why? At particle level the answer is beautifully simple.

Recall Chapter 6: weighing evidence such as definite and multiple proportions forced Dalton's atomic theory into being. Chemical reactions neither create nor destroy atoms nor change one element into another. That requires altering nuclei, Chapter 8's territory, beyond ordinary chemical energies. Atoms do only one thing: \textbf{recombine}.

Consider hydrogen burning in oxygen. Initially \chem{H_2} and \chem{O_2} each pair two atoms. Old partnerships break, absorbing energy, and atoms recombine into \chem{H_2O}, forming new bonds and releasing energy. Next chapter balances that energy account. Throughout, hydrogen remains hydrogen and oxygen oxygen, with none added or lost.

Mass conservation is therefore atom conservation reflected on a scale. Unchanged atom counts and individual masses necessarily give unchanged total mass. Conversely, if an experiment truly changes total mass, either something crossed your accounting boundary or you encountered a nuclear reaction. Chemists check boundaries first.

Magnify recombination into a dance changing partners: pairs release hands, move through the crowd, and join anew. Head count stays constant; partnerships change. Burning, rusting, bubbling, and color changes all hide the same event. Why this pairing instead of another---which lowers energy and how quickly switching happens---belongs to later chapters. Here, first establish the ledger: count everyone correctly.

\section{Balancing: Equal Counts of Every Atom}

We can write unchanged atom counts as a rule. Start hydrogen combustion with formulas:

\[\chem{H_2} + \chem{O_2} \rightarrow \chem{H_2O}\]

Count: 2 hydrogens and 2 oxygens left, but 2 hydrogens and only 1 oxygen right. Unbalanced. What now?

One discipline is essential: \textbf{change only coefficients before formulas, never subscripts within them}. Subscripts specify the molecular recipe. Turning \chem{H_2O} into \chem{H_2O_2} replaces water with hydrogen peroxide, switching goods merely to balance the accounts.

Instead adjust quantities: two hydrogen molecules produce two waters, using two oxygens:

\[\chem{2H_2} + \chem{O_2} \rightarrow \chem{2H_2O}\]

Now each side has 4 hydrogens and 2 oxygens. This is \textbf{balancing}, not a numerical game or making numbers look tidy, but translating constant atom counts into symbols.

Try slightly more complex methane combustion, the main natural-gas reaction:

\[\chem{CH_4} + \chem{O_2} \rightarrow \chem{CO_2} + \chem{H_2O}\]

Carbon already balances, one each. Four hydrogens left and two right require coefficient 2 on water. That makes four oxygens right, requiring 2 on oxygen:

\[\chem{CH_4} + \chem{2O_2} \rightarrow \chem{CO_2} + \chem{2H_2O}\]

Practice builds fluency, but the principle stays: count atoms, adjust quantities, and reduce coefficients to the smallest whole-number ratio. Coefficients give \textbf{particle-number ratios}: one methane consumes two oxygens. Chapter 5's mole bridge converts these into mass ratios for chemists' recipes, as we will calculate shortly.

A conditioned example: hydrogen peroxide solution decomposes slowly on standing. Add a little manganese dioxide powder and it bubbles vigorously:

\[\chem{2H_2O_2(aq)} \xrightarrow{\text{catalyst}} \chem{2H_2O(l)} + \chem{O_2(g)}\]

Balance it yourself: four hydrogens and four oxygens left; two waters use all four hydrogens and two oxygens, leaving one oxygen molecule. The catalyst above the arrow is not counted because manganese dioxide is not consumed overall; it speeds the route, as Chapter 30 explains, without needing a ledger entry. Laboratories often produce oxygen this way; the gas relights a glowing splint.

\begin{wuqu}
``Balancing just tidies numbers; if it fails, change a formula.'' This is the most dangerous shortcut. Writing hydrogen combustion's product as \chem{H_2O_2} balances oxygen but substitutes a completely different substance: peroxide can bleach and disinfect, while water is your daily drink. Each formula represents a real substance, not an erasable draft. Failure means your quantities or chosen reactants/products are wrong. Check the goods, not rewrite the books.
\end{wuqu}

\section{Everything One Equation Tells You}

A complete equation gives much more than who becomes whom:

\begin{itemize}[itemsep=2pt]
\item \textbf{Reactants} stand left of the arrow and \textbf{products} right; read the arrow as ``reacts to form.''
\item \textbf{Coefficients} give particle-count ratios, equivalently mole ratios, balanced through atom conservation.
\item \textbf{State symbols} follow formulas: solid (s), liquid (l), gas (g), and aqueous (aq), meaning water-containing. Chapter 25 showed dissolved particles differ from pure substances. State changes behavior dramatically: hydrogen gas flashes explosively at a spark, while dissolved hydrogen sits quietly.
\item \textbf{Conditions} appear above or below arrows. Heating often uses $\Delta$; catalysts are written above. A catalyst appears only in transit, not the overall account, offering an easier route without joining the final materials list (Chapter 30).
\item \textbf{Reversible arrows}, $\rightleftharpoons$, mark reactions running both ways simultaneously. Which direction dominates and how conditions shift it belongs to Chapter 31's equilibrium. For now, recognize the symbol.
\end{itemize}

A fully specified example is limestone calcination into quicklime, industrially around $\ud{900}{^{\circ}C}$:

\[\chem{CaCO_3(s)} \xrightarrow{\text{high temperature}} \chem{CaO(s)} + \chem{CO_2(g)}\]

This says solid calcium carbonate decomposes at high temperature into solid calcium oxide and carbon dioxide gas, in a $1:1:1$ particle ratio. Escaping kiln gas leaves lighter solid, but the books balance: some goods exited the accounting boundary through the chimney.

\begin{zhengju}
Mass conservation was not proclaimed by a sage; a balance forced it into view. Eighteenth-century chemistry was dominated by phlogiston: combustible matter supposedly contained a substance that escaped when burning. This explained lighter burned wood but struggled with heavier calcined metals, prompting advocates to give phlogiston negative mass. A theory surviving on patches is often nearing collapse.

In the 1770s, Lavoisier built experiments around \textbf{weighing}. Famously, he heated tin in sealed glass. Gray-white calx formed, but total sealed mass stayed unchanged. Opening let air hiss inward, showing some had been consumed. Reweighing showed the metal's gain exactly matched the missing gas's mass, down to the last fraction.

The alternative explanation collapsed: escaping phlogiston should have reduced sealed-vessel mass, yet did not. Metal taking something from air explained both mass gain and the incoming rush. Lavoisier's more precise mercury experiments measured the consumed fraction as about one-fifth, later called oxygen (Chapter 14).

Remember two lessons. His secret weapon was not brilliant intuition but meticulous balance use, with a sealed boundary and every mass change traced to a destination. Also, conservation began as an experimental generalization before anyone saw atoms. Only over a century later, with atomic theory established (Chapter 6), did it gain its particle explanation: matter persists because atoms do not change identity. Good laws first prove useful, then gain explanations that strengthen them.
\end{zhengju}

\section{Recipes from the Ledger: Stoichiometry}

A balanced equation is a \textbf{recipe}, not just a narrative. Coefficient ratios equal particle ratios equal mole ratios. Chapter 5's bridge brings masses into reach.

\begin{qiaoliang}
For $\chem{2H_2} + \chem{O_2} \rightarrow \chem{2H_2O}$, two portions \chem{H_2} plus one \chem{O_2} yield two \chem{H_2O}. In moles, \ud{2}{mol} hydrogen plus \ud{1}{mol} oxygen give \ud{2}{mol} water. In mass, using Chapter 5's relative molecular masses: hydrogen $2\times 2=\ud{4}{g}$ plus oxygen $\ud{32}{g}$ produce water $2\times 18=\ud{36}{g}$. Check: $4+32=36$, perfectly balanced. Rocket engineers use precisely such accounts, loading liquid hydrogen and oxygen in a $1:8$ mass ratio read from the equation rather than guessed.
\end{qiaoliang}

Real recipes rarely match exactly. The first ingredient exhausted determines where the reaction stops: the \textbf{limiting reactant}.

Start with sandwiches, each needing two bread slices and one cheese slice. Ten bread slices and four cheese slices make four sandwiches, using eight bread and leaving two. Cheese limits the yield; bread is in excess.

Chemistry is identical. Ignite \ud{4}{g} hydrogen (\ud{2}{mol}) with \ud{16}{g} oxygen (\ud{0.5}{mol}). The $2:1$ ratio requires \ud{1}{mol} oxygen for all hydrogen, but only \ud{0.5}{mol} is available, enough for \ud{1}{mol} hydrogen. Oxygen runs out, producing \ud{1}{mol} water (\ud{18}{g}) and leaving \ud{1}{mol} hydrogen. However long confined, that remainder cannot make water alone: you cannot cook without ingredients.

Now calculate the balloon experiment:

\[\chem{NaHCO_3} + \chem{CH_3COOH} \rightarrow \chem{CH_3COONa} + \chem{H_2O} + \chem{CO_2}\]

All coefficients are one, so each baking-soda unit produces one carbon dioxide. With molar mass about \ud{84}{g/mol}, two spoonfuls, roughly \ud{5}{g}, give about \ud{0.06}{mol} and theoretically the same carbon dioxide. At ordinary conditions \ud{1}{mol} gas occupies about \ud{24}{L}, so $0.06\times 24\approx\ud{1.4}{L}$, inflating a ball a little over ten centimeters across. Next time, calculate before experimenting and compare size with prediction. This habit protects chemists from mistakes: a basketball-sized balloon or no inflation immediately signals trouble.

Hospitals calculate too. Doctors often prescribe milligrams per kilogram because drug molecules interact with bodily substances in particle ratios. Dialysis and infusion ion concentrations likewise use precise proportions; an order-of-magnitude error can cause an accident. Stoichiometry is not an exam trick, but a universal language of proportionate action.

The maximum product calculated from reactant quantities is the \textbf{theoretical yield}. Ask a factory or laboratory whether anyone truly achieves 100\%: almost never. Actual divided by theoretical production is the \textbf{yield}. Why the discount? Several reasons lie beyond the ledger:

\begin{itemize}[itemsep=2pt]
\item \textbf{Side reactions}: materials may take another route to unwanted products; hot carbon may produce \chem{CO_2} or carbon monoxide, \chem{CO}.
\item \textbf{Incomplete reaction}: reversible reactions stall as opposing directions compete, leaving unconverted reactants. Chapter 31 calls this equilibrium.
\item \textbf{Losses}: films stay on vessels during transfer, filter paper retains product, and mother liquor holds dissolved crystals. Every operation leaks a little.
\item \textbf{Impure reactants}: Chapter 3 showed purity is difficult. Of \ud{10}{g} weighed, perhaps only \ud{9.5}{g} actually reacts.
\end{itemize}

Yield is thus a chemist's performance indicator. Raising it from 60\% to 90\% gives half as much product again from the same inputs and half the waste. Green chemistry's careful economics begins here.

Industry takes stoichiometry to extremes. A life-or-death example is the car airbag. A collision's electrical signal triggers sodium azide, \chem{NaN_3}:

\[\chem{2NaN_3(s)} \rightarrow \chem{2Na(s)} + \chem{3N_2(g)}\]

The design aims for about \ud{70}{L} nitrogen in 0.03 seconds. At room temperature, \ud{70}{L} is roughly \ud{3}{mol}; the $2:3$ ratio requires \ud{2}{mol} \chem{NaN_3}, about \ud{130}{g}. Too much overpressurizes and injures; too little fails to cushion. The dose comes step by step from the equation. Other ingredients immediately capture dangerous reactive sodium into harmless substances, also according to the ledger.

\section{What This Ledger Leaves Out}

Equations are useful, with equally clear limits. Learn what they \textbf{do not record}:

\begin{itemize}[itemsep=2pt]
\item \textbf{Energy}: hydrogen and oxygen make water, but the equation says neither how much heat nor why ignition is needed. Chapters 27 and 28 have that ledger.
\item \textbf{Speed}: even perfect balancing cannot tell one second from a thousand years. Paper thermodynamically ought to burn in air, yet may sit for decades. Chapter 29 handles rates.
\item \textbf{Pathway}: invisible intermediate steps may lie between start and finish. Equations record opening and closing balances, combining intermediate transactions. Catalysts change those transactions, not endpoints (Chapter 30).
\item \textbf{Guaranteed occurrence}: an equation gives permitted proportions, not a promise. Energy and equilibrium determine whether, in which direction, and how far reactions proceed (Chapters 28 and 31).
\end{itemize}

Equations govern how much, not whether, how fast, or by what route. Treating them as universal prophecy is a common misuse of symbols.

\section{Look Again at the Candle}

Where did the wax go? Its hydrocarbons, perhaps alkanes with twenty-some carbons, recombine with oxygen. Carbon enters carbon dioxide and hydrogen water vapor, dispersing into surrounding air with no atoms lost, only new partners. Treat candle, flame, and air as one account: wax plus oxygen becomes carbon dioxide plus water vapor in matching total grams. Disappearance merely means the boundary was too narrow.

The nail gains because oxygen enters the account; the balloon expands without mass change because gas never leaves. Atom conservation, balancing, coefficient ratios, limiting reactants, and theoretical yield now let you establish any reaction's ledger.

Look further: the wax's carbon drifts as carbon dioxide, perhaps entering a leaf next year as glucose and someone's bread years later. Chapter 16 showed elements cycling without spontaneous creation or disappearance. Each reaction ledger is a tiny entry in that immense cycle. In Volume Three, life itself becomes a meticulous accountant: cellular respiration and leaf photosynthesis count atoms one by one through balanced equations.

\section{Beyond the Ledger Lies Energy}

You may still wonder: if hydrogen and oxygen can react once balanced, why must ignition trigger the explosion? Why does an unlit candle not burn? Breaking costs, forming returns---how is this energy account kept?

The next ledger page records energy flows rather than atom counts. Chapter 27 balances it.

\begin{tiaozhan}
Balancing is algebra. For ethane combustion, set $a\chem{C_2H_6} + b\chem{O_2} \rightarrow c\chem{CO_2} + d\chem{H_2O}$. Conservation gives carbon $2a=c$, hydrogen $6a=2d$, and oxygen $2b=2c+d$. Let $a=1$: then $c=2$, $d=3$, $b=(2c+d)/2=7/2$. Fractions are fine; multiply all by 2: $\chem{2C_2H_6} + \chem{7O_2} \rightarrow \chem{4CO_2} + \chem{6H_2O}$. More generally, complete combustion of \chem{C_xH_yO_z} follows
\[\chem{C_xH_yO_z} + \left(x+\tfrac{y}{4}-\tfrac{z}{2}\right)\chem{O_2} \rightarrow x\chem{CO_2} + \tfrac{y}{2}\chem{H_2O}\]
Count atoms to verify. Unknown coefficients, conservation equations, and whole-number scaling balance most reactions. Electron transfer requires an additional electron-conservation rule, supplied in Chapter 33. Skipping this does not affect the main thread.
\end{tiaozhan}

\section*{Try It}

\begin{sikaoti}
\item Burning iron in oxygen forms \chem{Fe_3O_4}. Write an unbalanced equation, balance by counting, then calculate oxygen required and product formed from \ud{168}{g} iron. (Relative atomic masses: Fe $=56$, O $=16$.)
\item Draw differently colored circles for hydrogen and oxygen before and after $2\chem{H_2} + \chem{O_2} \rightarrow \chem{2H_2O}$. Count each atom type and note that circles are representations, not hard spherical atoms.
\item Each sandwich now needs two bread slices, one cheese, and one ham. With 12 bread, 5 cheese, and 7 ham, how many can you make, what limits production, and what remains? Compare this with question 4.
\item For $\chem{Zn(s)} + \chem{2HCl(aq)} \rightarrow \chem{ZnCl_2(aq)} + \chem{H_2(g)}$, combine \ud{6.5}{g} zinc (\ud{0.1}{mol}) with acid containing \ud{0.1}{mol} \chem{HCl}. What is theoretical hydrogen production, which reactant limits it, and what remains?
\item Suppose question 4 yields only 90\% of theoretical collected hydrogen. Give at least three possible causes, classifying them as side reactions, incomplete reaction, handling loss, or impurity.
\item Ammonium bicarbonate fertilizer vanishes readily on heating: $\chem{NH_4HCO_3(s)} \xrightarrow{\text{heating}} \chem{NH_3(g)} + \chem{H_2O(g)} + \chem{CO_2(g)}$. A farmer finds an open bag has become lighter and thinks mass conservation failed. Explain using accounting boundaries and design an experiment demonstrating conservation.
\end{sikaoti}
